% `tex.web` §379 and §367: an `\input` that turns up while a file name is
% being scanned is put back, not obeyed.
%
%     input: if chr_code>0 then force_eof:=true
%       else if name_in_progress then insert_relax
%       else start_input;
%
%     procedure insert_relax;
%     begin cur_tok:=cs_token_flag+cur_cs; back_input;
%     cur_tok:=cs_token_flag+frozen_relax; back_input; token_type:=inserted;
%     end;
%
% §526 raises `name_in_progress` over the name itself and §1258 raises it
% again over the `at`/`scaled` keyword that may follow, so it catches an
% `\input` on either side of a font's size.
%
% We ended the name at the control sequence and backed it up unexpanded. The
% file still opened in the right place and the name still came out right — but
% the inserted `\relax` is a command, and §367's `tracing_commands>1` traces
% the expansion of `\input` as well, so TeX writes two lines here that we
% wrote none of.
%
% `\input` is not expandable in this engine — `main_control` dispatches it —
% so the flag is consulted at the one place TeX ever raises it, rather than in
% `expand` where TeX keeps it.
\catcode`\{=1 \catcode`\}=2
\tracingonline=1 \tracingcommands=2
\errorcontextlines=5

\immediate\openout0=isub.tex
\immediate\write0{\string\message{PROBE:sub}\string\showthe\string\escapechar}
\immediate\closeout0

% The name runs straight into the `\input`.
\font\a=cmr10\input isub

% ...and here the size scan is what meets it: §1258's half of the flag.
\font\b=cmr10 at 12pt \input isub

\showthe\escapechar
\end
